\documentclass{article}
\usepackage{hyperref}
\usepackage{Style}

\nocite{*} % Comentar si quiero citar
%\addbibresource{bibliografia.bib} % Quitar el comentado si quiero usar bibliografia

\begin{document}

\begin{minipage}{2.5cm}
    \includegraphics[width=2cm]{imagen_puc.jpg}
\end{minipage}
\begin{minipage}{12cm}
    {\sc Pontificia Universidad Católica de Chile\\
    Facultad de Matemáticas\\
    Departamento de Matemática\\
    Profesor: Manuel Cabezas -- Ayudante: Benjamín Mateluna}
\end{minipage}
\vspace{1ex}

{\centerline{\bf Introducción a la Geometría - MAT1304}
\centerline{\bf Ayudantía 3}}
\centerline{\bf 17 de marzo de 2026}

\vspace{1cm}
\noindent\textbf{Problema 1.} Sea $\triangle ABC$ equilátero. Sea $P$ un punto en el interior del
triángulo y suponga que $D,F,E$ son puntos de $\overline{AB},\overline{AC},\overline{BC}$ tales 
que $\overline{DP},\overline{EP}$ y $\overline{FP}$ son perpendiculares a 
$\overline{AB},\overline{AC},\overline{BC}$ respectivamente. Demuestre que 
$\overline{DP}+\overline{EP}+\overline{FP}$ es la altura del triángulo equilátero.

\vspace{5mm}
\noindent\textbf{Problema 2.} Dado un triángulo $\triangle ABC$ denotamos por $M$ el punto medio
de $\overline{AB}$. Demuestre que si $\angle BCA=90^{\circ}$, entonces 
$\overline{AM}=\overline{BM}=\overline{CM}$.

\vspace{5mm}
\noindent\textbf{Problema 3.} Considere las medidas $a,b,c,d,e,f$ de la figura de abajo. Asuma que
son enteros positivos. ¿Puede ser que $a,c,e$ sean pares y $b,d,f$ sean impares? Justificar. (Los
segmentos son concurrentes)

\begin{center}
    \includegraphics[scale=0.25]{P3.png}
\end{center}

\noindent\textbf{Problema 4.} Sea $\triangle ABC$ un triángulo cualquiera. Consideremos los puntos 
$D,E,F$ en $\overline{BC},\overline{AC},\overline{AB}$ respectivamente, tales que 
$\overline{AD},\overline{BE}$ y $\overline{CF}$ concurren dentro del triángulo. Muestre que 
si $\overline{FE}$ y $\overline{BC}$ son paralelas, entonces $\overline{BD}=\overline{DC}$.

\vspace{1cm}
\noindent\textbf{\underline{Ejercicios Propuestos:}}

\vspace{5mm}
\noindent\textbf{Extra 1.} Demuestre que todo triángulo se puede dividir en cuatro triangulos 
isósceles. \textbf{Hint:} Usar el segundo problema.

\vspace{5mm}
\noindent\textbf{Extra 2.} En la figura tenemos un triángulo equilátero de área $\sqrt{3}$ dentro
de un cuadrado

\begin{center}
    \includegraphics[scale=0.25]{E2.png}
\end{center}

\begin{enumerate}
    \item Encuentre la medida del lado del triángulo equilátero.
    \item Muestre que dos de los cuatro triángulos en la figura son congruentes.
    \item Calcule el área del cuadrado.
\end{enumerate}

%\printbibliography % Quitar el comentado si quiero usar bibliografia

\end{document}